% ---------------------------------------------------------------------------
% Setup additions for Section "The SHIFT-SURV Benchmark"
%
% Two items from the review:
%   (1) "metriche che usi"  -> the metrics are never defined in the draft.
%       New subsection below; put it after "Methods evaluated", i.e. immediately
%       before "Protocol", so the reader has methods + metrics before Results.
%   (2) "invariant oracle -> oracolo di learning, piu' vicino alla realta'"
%       -> replacement for the existing "Invariant oracle" subsection.
%
% The shift parameters the review also asked for are in table_scenarios.tex
% (generated from configs/generator.yaml), which replaces the current
% tab:scenarios environment.
% ---------------------------------------------------------------------------


% ===== REPLACES the current \subsection{Invariant oracle} =====================

\subsection{Invariant oracle}
Because the true latent factors are known, we fit a logistic regression on
$u_{\mathrm{inv}}$ under the same training and evaluation protocol as every
other method. The oracle is deliberately a \emph{learning} oracle rather than an
information-theoretic bound: it is a model of the same family, fit on the same
finite sample, with the same train/validation split and the same horizon, and
differs from the compared methods in one respect only --- it reads the invariant
representation instead of the noisy observation $x$. It is therefore not the
best achievable predictor, but the best that a learner of this class could reach
if the transportable part of the signal were handed to it directly, which makes
the headroom we report a conservative and attainable target rather than an
unreachable ceiling. The oracle is not deployable, since $u_{\mathrm{inv}}$ is
unobserved. Its discrimination sets the reference against which we quantify how
much of the cross-center gap each method actually recovers, and its calibration
provides a floor on the calibration error achievable under a given shift.


% ===== NEW subsection, insert before \subsection{Protocol} ===================

\subsection{Evaluation metrics}
Every method outputs a predicted probability that the event occurs by the
horizon, so we evaluate two properties that a transported risk model must have
and that can fail independently. Discrimination is measured by the area under
the receiver operating characteristic curve (AUROC), which depends only on how
patients are ranked and is unchanged by any monotone rescaling of the
predictions. Calibration is measured by the expected calibration error (ECE)
over $B=10$ equal-width probability bins,
\begin{equation}
\mathrm{ECE} \;=\; \sum_{b=1}^{B} \frac{n_b}{n}\,
  \bigl\lvert \bar{y}_b - \bar{\hat{p}}_b \bigr\rvert ,
\end{equation}
where $n_b$ is the number of target patients whose predicted probability falls
in bin $b$, and $\bar{y}_b$ and $\bar{\hat{p}}_b$ are the observed event rate and
the mean predicted probability within that bin. Reporting both is necessary
rather than decorative: prior adjustment and Platt scaling are monotone
transformations of the ERM scores and so leave AUROC unchanged by construction,
and an evaluation restricted to discrimination would be blind to precisely the
methods that act on calibration.

Both metrics are computed on the held-out target center, whose labels are used
only for evaluation and never for fitting or for hyperparameter selection. We
fix in advance a margin of $0.02$ AUROC as the smallest improvement that would
be clinically meaningful, and read a confidence interval contained entirely
within $\pm 0.02$ as evidence of no meaningful gain, rather than as an
inconclusive comparison.
